\chapter{如何說話}

這場講座把說話視為一門有結構、有技巧、也有後果的技藝。它並不是從訣竅講起，而是先從迫切性切入，再把這種迫切性收束為一個關於溝通品質的公式，之後才逐步展開整套機制：如何開場、如何不讓聽眾的注意力飄散、如何使用黑板與道具、為什麼投影片往往失靈、求職演講幾乎一開始就必須建立什麼，以及講者應當如何收尾。這樣的次序本身就是教學設計的一部分。我們先被說服：這裡確實有東西可學；接著才看見，這些知識在實際運作中究竟是什麼樣子。

\section{開場鉤子、溝通公式與承諾}

Winston 以一個強而有力的類比開場。他說，《軍事司法統一法典》會懲罰那種讓士兵手無寸鐵就上戰場的軍官。對學生也應當有類似的保障，因為學生不該在沒有溝通能力的情況下被送入人生。重點不在文學性的修飾，而在於：思想不會自己傳播。

接著他用幾乎不能更直白的方式給出排序：

\begin{equation}
\text{說話} \succ \text{寫作} \succ \text{思想的品質}.
\end{equation}

這裡的 \(\succ\) 只是編者用來對應他那句 ``in that order'' 的簡寫。這個主張並不是說，思想在本質上比表達更不重要，而是說，一個人在世界中的成敗，往往會先被這些因素各自失敗或成功的先後次序所篩選。

接著，講座把焦點收束到第一個正式對象。若依據可靠的逐字稿，而不是稍後黑板上殘留的痕跡，我們可以把開場公式謹慎地寫成標準形式

\begin{equation}
Q = f(K,P,T),
\end{equation}

其中 \(Q\) 表示說話或寫作的品質，\(K\) 是知識，\(P\) 是練習，\(T\) 是天賦。

\begin{remark}
在之後那張出現自行車輪的教室畫面中，這個公式仍只留下部分背景殘跡，大致像是 \(\text{Quality}=f(K,P,\ldots)\)。因此，這裡的完整三變數形式是根據逐字稿重建，而不是單靠截圖完整轉錄而來。
\end{remark}

講座隨即告訴我們該如何理解這個公式。它的各個變數並不對稱。Winston 說，\(T\) ``very small''。若用簡潔的數學速記來忠實概括他的意思，而不是照抄他黑板上的記號，我們可以寫成

\begin{equation}
T \ll K,P.
\end{equation}

關鍵詞不只是 ``practice''，更是 ``knowledge''。Winston 說得非常明白：真正重要的是你知道什麼。

接著，他不是用抽象論證，而是用一個故事來檢驗這個公式。Mary Lou Retton 是一位非凡的運動員，但她是滑雪新手。Winston 在運動天賦上遠不如她，卻是更好的滑雪者。為什麼？因為在這個問題上，他擁有相關的 \(K\) 與 \(P\)，而她主要擁有的是 \(T\)。這個軼事不是裝飾公式的點綴，而是第一個證據，說明這個公式必須被認真看待。

\subsection{問題與回答}

\paragraph{問題。} 為什麼 \(T\) 很小？

\paragraph{回答。} 因為單靠天賦本身，還不足以構成特定任務的勝任能力。滑雪的例子把這一點說得再鮮明不過。Retton 擁有驚人的一般運動天賦，但 Winston 擁有滑雪的知識與滑雪的練習。因此在雪板上，他更強。這場講座要我們把這個結構直接轉移到說話上。天生資質中等、但有良好典範、習慣與準備的人，可以勝過一位天分更高、卻訓練不足的講者。

從這個軼事出發，講座立刻轉向承諾。若這個公式是對的，那麼改進就不是神祕之事。聽眾被承諾將獲得一整套說話技巧的武器庫，而 Winston 又補上一個重要但書：這個過程不是線性的。也許其中某一項技巧，甚至只是一項，就會成為那個幫某人拿到工作的關鍵。這不是一句隨口帶過的話，而是從理論通往整堂課其餘內容的動機橋樑。

\section{參與規則與如何開場}

Winston 直到做出這個承諾之後，才開始施加紀律。不准開筆電、不准用手機，這條規矩並不是作為禮儀問題提出的，而是作為一個認知事實提出的。他說，人只有一個語言處理器。如果那個處理器正在讀電子郵件或瀏覽網頁，它就沒有在聽。更糟的是，損害還具有社會性：周圍的聽眾會被分散注意力，而當打開的筆電持續映入眼簾時，講者的表現也會變差。

這是整場講座的核心觀念之一。它稍後還會再次出現，成為反對文字過密投影片的最深層理由。而在此處，它執行的是另一項功能：在正式講述開始之前，先把通道清乾淨。

下一步同樣是刻意安排的。Winston 說，``First thing to talk about,'' ``is how to start.'' 他立刻否定一條耳熟能詳的建議：以笑話開場。理由不是笑話不配進入嚴肅的演講，而是時機問題。在演講剛開始時，整個房間還在適應講者的聲音、速度與節奏。聽眾尚未與講者同步。

\subsection{問題與回答}

\paragraph{問題。} 為什麼不要以笑話開場？

\paragraph{回答。} 因為笑話需要一群已經安定下來的聽眾，而一場演講的開頭給我們的，恰恰是一群尚未安定的人。人們還在進場、還在收拾東西、還在調整自己對講者聲音參數的感受。因此，笑話往往出現在聽眾尚未準備好接收它的時刻。開頭真正應該放上的，是一個賦能承諾：清楚說明在這一小時結束時，聽眾將會知道什麼，而那是他們在一開始尚且不知道的。

這是整場講座最漂亮的自反特徵之一。Winston 不只是建議要提出賦能承諾；他其實已經在這麼做。講座先以迫切性開場，再提出公式，再給出一個帶有證明意味的軼事，而現在又承諾會提供可用的技巧。這場演講本身，已經被拿來充當它自己的範例。

\section{四條啟發法}

一旦聽眾被說服：這裡確實有真知識可學，Winston 就提出四條他說話時時刻放在心上的示範性啟發法：繞著主題循環、為觀念築起柵欄、使用語言標點，以及提出問題。

\paragraph{循環。} 第一條啟發法，是圍繞主題講一遍，然後再繞一遍。Winston 提醒我們那句老建議：先告訴他們你要告訴他們什麼，接著真的告訴他們，然後第三次再告訴他們一次。他的理由不是因為聽眾愚鈍，而是因為聽眾的注意力本來就是間歇性的。他說，不管講的是什麼，在任何一個時刻，教室裡大概總有 \(20\%\) 的人處於神遊霧化狀態。

這條啟發法自然引出一個小小的編者模型。若 \(p_{\mathrm{fog}}\) 表示一位聽者錯過某一次講述的機率，那麼 Winston 那個 \(20\%\) 的估計便給出

\begin{align}
p_{\mathrm{fog}} &\approx 0.2,\\
\mathrm{Prob}(\text{三次都錯過}) &\approx p_{\mathrm{fog}}^3 = 0.2^3 = 0.008,\\
\mathrm{Prob}(\text{至少接住一次}) &\approx 1-p_{\mathrm{fog}}^3 = 0.992.
\end{align}

\begin{remark}
這是我們的說明性詮釋，不是 Winston 自己的計算。講座真正的重點是實際層面的，而非機率層面的：重複會創造多次重新進入的機會。
\end{remark}

\subsection{問題與回答}

\paragraph{問題。} 為什麼要說三次？

\paragraph{回答。} 因為說話是在時間中發生的，而聽眾不會完美地靜止不動。只說一遍是脆弱的。第二遍與第三遍，尤其是從稍微不同的側面接近同一個重點時，會讓溝通變得穩健。因此，``循環'' 不是空洞的重複，而是對注意力漂移所做出的工程性回應。

\paragraph{築起柵欄。} 第二條啟發法，是透過對比來下定義。若我們希望聽眾知道我們的觀念是什麼，也應幫助他們看見它不是什麼。Winston 的簡單例子是拱：這是一個拱；這不是拱；那也不是拱。技術性的形式同樣有用：我的演算法與 Jones 的演算法相似，不同之處在於，他的是指數級，我的是線性。當我們替一個概念築起柵欄時，它的邊界也就變得銳利起來。

\paragraph{語言標點。} 第三條啟發法，來自對注意力的同一診斷。若人們會從公車上掉下來，那麼演講中就需要有一些接縫，好讓他們重新上車。明確的列點、回顧與轉場語句，就是這些接縫。Winston 自己的講座正示範了這種方法：他從第一個觀念轉到第二個，接著又明白宣布現在要進入第三個。

\paragraph{提出問題。} 第四條啟發法，是最直接的重新喚回注意力的裝置。一個問題會打斷被動接收，迫使整個房間把自己重新對準某個問題。但即便如此，講座仍補上一條有分寸的規則：你必須等得夠久。

\begin{equation}
t_{\mathrm{wait}} \approx 7\,\mathrm{s}.
\end{equation}

對講者來說，七秒彷彿漫長得像永恆。也正因如此，這條規則才有必要。整個房間需要思考時間。同時，問題本身也必須謹慎挑選：不能簡單到讓人回答起來覺得尷尬，也不能困難到誰都說不出話。

在這一節的末尾，Winston 停下來說了一句很耐人尋味的話。如果聽眾此刻已經相信，關於說話確實 ``something to know''，那麼他就已經成功了。這正是整場講座轉折的鉸鏈。直到此刻，他才把視野擴大。

\section{時間與場所作為論證的一部分}

Winston 說，議程上的下一件事，是時間與場所。這個轉折說得非常明白，而講座也希望我們感受到：這樣的擴展是先前鋪墊出來的。先前我們處理的是一場演講的內在機械結構；現在則要考慮，這些機制是在什麼條件下才能運作。

理想的講課時間是上午 \(11\) 點。Winston 給出的理由既實際又與身體狀態有關：人已經清醒，不是在飯後，也還沒有開始陷入疲憊。理想的教室應當光線充足，因為昏暗的房間會發出睡意的訊號。它也應當事先勘察過，也就是說，預先查看，以便把臨場意外降到最低。最後，房間裡還應有相當數量的人。一個空了一半的大廳，會自行製造出令人洩氣的社會劇場。

這些要點在講座最好的一張黑板畫面中仍然以視覺形式保存下來。那張截圖之所以珍貴，不是因為它包含了什麼推導，而是因為它保留了黑板的架構。

\begin{figure}[h]
\centering
\includegraphics[width=0.95\textwidth]{lecture_11_figure_03.png}
\end{figure}

左側標題是 \emph{範例}，下面列出四條啟發法。``循環'' 旁邊那個環狀圖形不是正式圖解，但它確實以視覺方式編碼了反覆多次講述的觀念。右側標題是 \emph{時間 \& 場所}，中間以一道長長的波浪線分開。這塊黑板本身已經在做 Winston 所說一塊好黑板該做的事：它是在空間中組織一場演講，而不只是用語言組織它。

把它重新整理後，這樣的組織會更容易讀懂。右半邊的確切字句，單靠畫面本身其實並不完全可靠，因此下面的項目是參照逐字稿補全的，同時保留黑板整體的幾何布局。

\begin{figure}[h]
\centering
\begin{tikzpicture}[x=1cm,y=1cm]
\draw[thick] (0,0) rectangle (13.5,6.8);

\node[anchor=west] at (0.6,6.0) {\Large 範例};
\node[anchor=west] at (0.8,5.0) {循環};
\draw[thick] (3.0,5.0) ellipse (0.45 and 0.95);
\draw[thick] (3.45,4.95) ellipse (0.45 and 0.95);
\draw[->,thick] (4.0,5.0) -- (4.5,5.0);

\node[anchor=west] at (0.8,4.0) {築起柵欄};
\node[anchor=west] at (0.8,3.0) {語言標點};
\node[anchor=west] at (0.8,2.0) {提出問題};

\node[anchor=west] at (6.0,6.0) {\Large 時間 \& 場所};
\draw[thick] plot[smooth] coordinates {(5.8,4.85) (6.8,4.95) (7.8,4.8) (8.8,4.95) (9.8,4.75) (10.8,4.9) (11.8,4.7) (12.8,4.8)};
\node[anchor=west] at (6.5,4.1) {上午 11 點};
\node[anchor=west] at (6.5,3.3) {光線充足};
\node[anchor=west] at (6.5,2.5) {事先勘察};
\node[anchor=west] at (6.5,1.7) {有人氣};
\end{tikzpicture}
\end{figure}

這場講座真正要說的是：時間與場所不只是後勤細節，它們是注意力動力學中的變數。

\section{黑板與道具：為什麼實體媒介有效}

一直到談過時間與場所之後，Winston 才轉向工具。這個次序很重要。工具不是在真空中被選擇的；它是在某個房間裡、某個時刻、為了某個目的而被選擇的。

最基本的區分其實很簡單。若目的是告知、教學或講課，黑板就是合適的工具。投影片則更適合快速展示觀念，例如在會議演講與求職演講中。接著，Winston 提出三個理由，說明為什麼在真正的教學裡他偏好黑板。

第一，它具有圖像品質。黑板讓我們可以分組、分隔、畫底線、勾勒草圖。第二，它有一種速度屬性：書寫的速度，大致就是聽眾能吸收的速度。投影片則可能跑得比思考還快。第三，黑板替講者的雙手提供了一個可以指向的目標。這不只是一句玩笑。公開演說中的許多尷尬，都來自於身體忽然變得過度自覺；而黑板讓這個身體有一件誠實的工作可以做。

從這裡開始，講座把視野擴展到道具。主導原則依然是：先有例子，再有理論。Ibsen 在 \emph{Hedda Gabler} 裡的火爐之所以有效，是因為它就擺在那裡，逐漸積聚張力，慢慢變得不可避免。道具之所以令人難忘，不是因為它被描述了，而是因為它被置入場景、被演了出來。

同樣的邏輯也出現在教室裡。Papert 的自行車輪不是裝飾品，而是一個承載一整條推理線索的裝置。

\begin{figure}[h]
\centering
\includegraphics[width=0.78\textwidth]{lecture_11_figure_04.png}
\end{figure}

這張圖的重要性，首先在於那個身體性的示範：車輪豎直立著，輪轂握在講者手中，道具完整地呈現在視野裡。背景中仍可看見開場黑板公式的一部分殘跡，大致像是 \(\text{Quality}=f(K,P,\ldots)\)，以及更早先那個火爐草圖。但這張圖在數學上的負荷，落在道具本身，而不在那些殘留痕跡上。

\subsection{問題與回答}

\paragraph{問題。} 旋轉的輪子會往哪個方向走？

\paragraph{回答。} Winston 的意思是，死記下來的手則其實不是思考這個示範的好方法。他帶著修辭意味地說，許多機械工程師最後大約只有 \(50\%\) 的成功率，這等於表示那條規則已經退化成擲硬幣。更好的方法，是改變我們對問題的表徵方式。

\begin{enumerate}
\item 先讓輪子轉起來，再用膠帶在輪框上做一個小標記。
\item 暫時拒絕把整個輪子當成一個整體來想。
\item 看著那段做了標記的輪框越過最上方。
\item 對那一小段施加那股氣流，或等效的局部推力。
\item 問：這一小塊現在想往哪裡去？
\item 接著再看下一小塊抵達同樣的位置。
\item 同樣的局部推理再次適用。
\item 因此，整個輪子必須沿著對應的方向繼續前進。
\end{enumerate}

這裡的成就並不是完整的陀螺分析。講座從未把那套分析寫出來，我們也不該把它硬行補進來，彷彿它原本就在場。這裡真正的成就是教學上的。這個道具讓人得以使用一種局部小片段的論證，取代死記硬背的符號規則，轉而訴諸一幅更可靠的圖像。

第二個道具例子，是 Alan Lazarus 用擺球示範能量守恆。一顆沉重的球被提到鼻尖高度，由靜止釋放，先盪開，然後再盪回來，輕觸但不會撞毀示範者的鼻子。相關的方程式是那個經典形式：

\begin{equation}
E = K_{\mathrm{kin}} + V_{\mathrm{pot}} = \text{常數}.
\end{equation}

講座並沒有把它在黑板上推導出來，而且也不需要。重點在於：這個示範會把這條定律以身體性的方式刻進記憶裡。Winston 說，真正的危險並不是守恆失效，而是人總想推一把，而不是單純放手。

一直到火爐、車輪與擺球都講完之後，Winston 才提出一個一般性的說明。他在自己的 ``lunatic fringe'' 理論中說，粉筆與道具之所以有效，是因為同理性的鏡像作用。當我們看著一隻手書寫，或看著一顆球擺動時，我們幾乎會感到自己也在書寫、也在擺動。投影片或圖片並不會產生同樣的具身參與。講座再一次忠於自己的方法：例子在前，解釋在後。

\section{投影片、特殊情況與求職演講的門檻}

Winston 說，還有一項工具必須討論：投影片。與先前的區分，在黑板與道具都已經確立之後，如今以更強的力道再被重申一次。投影片是用來展示觀念的，不是用來教觀念的。

他的抨擊從一條幾乎不容例外的定律開始：人們的投影片永遠太多，而那些投影片上的字也永遠太多。接著，講座就變成一份投影片罪狀目錄。你不能把投影片上的內容整段讀出來。你不能站得離投影片那麼遠，以致聽眾被迫在螢幕與講者之間來回進行一場視覺網球。你應該剝除背景垃圾、多餘的 logo，以及那些講者本來就已經說出口的標題。

更深一層的論證，仍是開頭那個單一語言處理器的主張。字太密的投影片會把聽眾徵召去閱讀；而整個房間一旦開始閱讀，就不會在聽。Winston 又用一個實驗來強化這一點：當部分內容由口頭說出、另一部分則出現在投影片上時，學生對投影片內容的記憶反而更好。其中一位學生事後還說，講者一直講話其實很干擾。這個笑點之所以成立，正因為那個原理是真的。

接著，講座把建議說得更尖銳。一張投影片通常只應該承載少量文字。若其中確實有重要的歷史文本，講者就應明白給聽眾時間去讀。而這裡有一個有趣的例外：Winston 的 \emph{hapax legomenon} 投影片，那種密到幾乎不可能閱讀的投影片，只有在整個重點就是要展示複雜性本身時，才允許出現一次。但一次就是一次。

在字體排印上，規則是很實際的：

\begin{equation}
40\text{--}50\,\mathrm{pt}\ \text{是安全的}, \qquad 35\,\mathrm{pt}\ \text{就已經開始偏小了。}
\end{equation}

講座的重點不只是 \(35\) 點字可能難以辨認，而是講者通常只有在企圖把太多語言塞進投影片時，才會去摸索 \(35\) 點字。

Winston 還加上更多罪狀：雷射筆，因為它迫使講者轉身背對聽眾；過重的整套投影片，而這一點可以透過把投影片全部印出來、攤在桌上診斷出來；以及那種沒有呼吸感、沒有留白、也沒有圖像的投影片。這些警告累積起來的教訓是：投影片應該只是口頭推理線索的調味品，而不是牆上另行上演的一場競爭性講座。

從工具出發，講座接著轉入特殊情況。第一種是 informing，也就是說，正如 Winston 此刻正在做的那種說話。在這裡，賦能承諾又回來了，但強調點改變了：若要激發人心，光是把資訊組織好還不夠。你可以告訴學生他們做得到，可以幫助他們用新的方式看見一個問題，或者只是單純展現你對主題的真誠熱情。資源配置的故事就是他的例子。首先，我們被展示一個災難性的暴力計算；接著，我們被承諾，只要做一個小小的改動，而這個改動會在接下來五十分鐘內被說清楚，就能把一個本來需要整個太陽系壽命的計算，壓縮成幾秒鐘。這個承諾既有說明性，也令人振奮。

接下來是一段較小的岔題，談的是教師常說自己在做的事：教人如何思考。Winston 問出了這種說法自然會招來的下一個問題：怎麼教？他的回答是，我們是說故事的動物。因此，所謂教思考，就是提供故事、提出關於故事的問題、提供分析的機制、教人如何把故事拼接起來，以及如何評估它們的可靠性。這段岔題之所以重要，是因為它把呈現不只連到說服，也連到教育。

接著，口試又提供了另外兩條建議。考生必須把自己的工作放入脈絡之中，也就是說，把它放回更大的領域裡，並解釋這個問題為什麼重要。再者，練習必須是真正的練習，而不是在那些早已知道內容、會默默替你補完缺漏說明的人面前排演。

\subsection{問題與回答}

\paragraph{問題。} 在求職演講中，你有多少時間？

\paragraph{回答。} 大約五分鐘。

\begin{equation}
t_{\mathrm{job}} \approx 5\,\mathrm{min}.
\end{equation}

Winston 說，在那個時間窗之內，有兩件事必須建立起來。聽眾必須看見，這位候選人有 vision；同時也必須看見，這位候選人確實做出了一些東西。所謂 vision，指的是一個有人在乎的問題，加上一種在方法上真正新的做法。而 ``done something'' 則表示有具體進展，而不只是抱負。

講座透過具體建議把這件事落實下來。講者應當把解決這個問題所需的步驟逐條列出，指出其中哪些步驟其實已經完成，然後在最後用一張 contributions 投影片把那些步驟映照回來。於是，一個緊湊的技術演講結構就浮現出來：

\begin{enumerate}
\item 說明那個值得人關心的問題，
\item 指出方法中的新意是什麼，
\item 列出真正取得進展所必需的步驟，
\item 展示其中若干步驟其實已經完成，
\item 最後以一張對應前述清單的 contributions 投影片收束。
\end{enumerate}

這不只是禮節而已。這是一套理論，用來解釋陌生人如何在極短時間內判斷：眼前這位講者究竟擁有一個研究計畫，還只是握有一堆彼此孤立的事實。

\section{如何被記住，以及如何知道該停下來}

講座的最後一個樂章始於求職演講之後。Winston 說，就算工作拿到了，人仍然必須思考，如何讓自己的工作價值被辨認出來。他用 Julia Child 的軼事來鋪陳這一點。人們不斷走向她，對她說她改變了他們的人生。當被問到出名好不好玩時，她的回答很乾：人會習慣。Winston 的回應則更耐人尋味：人永遠不會習慣被忽視。這則故事的寓意是，包裝不是虛榮，而是讓觀念在這個世界上得到公平機會的一部分。

接著，他引入 Winston's Star，也就是讓工作變得令人難忘的五部結構：

\begin{itemize}
\item 符號，
\item 標語，
\item 驚奇，
\item 顯著觀念，
\item 故事。
\end{itemize}

那個古老的學拱例子正好把這顆星逐一展開。拱本身就是符號。``One-shot learning'' 是標語。驚奇之處在於：如果學習者夠聰明，一個例子就足以教會某件明確的東西。顯著觀念則是那個幾乎命中的例子：不僅僅是一個反例，而是一個近乎是例子的例子，它揭露出概念真正要求的是什麼。然後，所有這些元素都必須被綁進一個故事裡，說明這套系統如何運作，以及它為何重要。Winston 特別強調，這裡的 ``salient'' 並不是抽象地指 ``重要''，而是指那種足夠凸出、足以留在記憶中的觀念。

一直到談完記憶，講座才進入收束。這時，講座開始示範一些糟糕的結尾。不應該用一長串巨大的合作名單作結，因為就在聽眾判斷講者究竟做了什麼的那一刻，這會把講者縮小。不應該把最後一張投影片浪費在一片空白上。若真正重要的是貢獻，也不應該以 ``conclusions'' 作結。

\subsection{問題與回答}

\paragraph{問題。} 最後一張投影片與最後一句話應該做什麼？

\paragraph{回答。} 最後一張投影片應該展示貢獻，而最後一句話則應該完成一個真正的收束動作，而不是一個軟弱的禮貌姿態。

當人們提問時，當他們起身離開房間時，最後一張投影片仍會停留在那裡。因此，它應當把整份工作的最強壓縮版本持續留在視野之中。``Contributions'' 是 Winston 偏愛的標籤，正因為它映照了前面那份清單：哪些事情必須完成，以及哪些事情其實已經完成。

最後一句話則是另一門藝術。Winston 說，笑話放在結尾是可行的，因為到了那時，聽眾已經適應了講者，也能接住那個節奏。但他不建議把 ``thank you'' 當成主要的結語動作。這樣的結尾，彷彿是在暗示：聽眾留下來只是出於禮貌，而不是因為這場演講值得一聽。因此，他轉而去尋找其他範例：政治性的祝禱語、典禮性的解散公式、音樂會的慣例，以及最後他在這裡最想推薦的那個動作，也就是向聽眾與這個地方致意。

這一點同樣由講座本身親自演出。Winston 告訴聽眾，光是他們出現在這裡，就已經表明他們理解呈現是重要的，而他向他們致意。真正的收束動作完成之後，最後才補上一句有禮的致謝。這個順序很重要。

\section{總結}

這場講座以刻意的收束與擴展向前推進。它先收束為一個公式，
\[
Q=f(K,P,T),
\]
然後立刻藉由縮小 \(T\) 來打破這個公式的對稱性。接著它擴展到啟發法，再擴展到房間，再擴展到媒介的選擇，再擴展到口試與求職演講這類特殊情況，最後又重新收束為幾個緊湊的摘要裝置：Winston's Star、contributions 投影片，以及收尾的致意動作。

把整場講座統一起來的是一個單一而實際的問題：我們要如何讓觀念在與時間、分心、具身性與記憶接觸之後，仍有最大的存活機會？Winston 的答案不是某一個技巧，而是一整套 repertoire。我們要激發動機、做出承諾、反覆循環、築起柵欄、加上標點、提出問題、挑選房間、挑選媒介、在可能時用道具承載論證，並在最後留下對已完成之事最有力、最凝練的表述。